% TOP_PROBLEM: {"problemId": "problem.shub-smale-conjecture-for-integer-zeros", "canonicalTitle": "Shub-Smale τ-Conjecture for Integer Zeros", "releaseRank": 224, "primaryDomain": "theoretical_computer_science", "primaryDomainLabel": "Theoretical computer science", "displayStatus": "open", "releaseStatus": "open"}
% !TeX program = lualatex
% Definition and sources only; this file is not an LLM solution attempt.
\documentclass[11pt]{article}
\usepackage[margin=25mm]{geometry}
\usepackage{fontspec}
\setmainfont{Latin Modern Roman}
\usepackage{amsmath,amssymb,mathtools}
\usepackage{unicode-math}
\setmathfont{Latin Modern Math}
\usepackage{xurl,hyperref}
\hypersetup{colorlinks=true,urlcolor=blue}
\setcounter{secnumdepth}{0}
\setlength{\parindent}{0pt}
\setlength{\parskip}{5pt}
\setlength{\emergencystretch}{3em}
\begin{document}
\section*{\texorpdfstring{Shub-Smale \ensuremath{\tau}-Conjecture for Integer Zeros}{Shub-Smale τ-Conjecture for Integer Zeros}}\label{224}
\subsection{Definitions and mathematical statement}
A constant-free straight-line program begins with registers $1$ and $X$ and successively creates registers by addition, subtraction, or multiplication of earlier registers. Let $\tau(f)$ be the minimum number of these operations needed to output $f\in{\mathbb{Z}}[X]$. For $f\ne0$, let $Z_{{\mathbb{Z}}}(f)=\{a\in{\mathbb{Z}}:f(a)=0\}$. The conjecture asks for absolute $a>0,c\ge1$ with
\[
 |Z_{{\mathbb{Z}}}(f)|\le a(1+\tau(f))^c.
\]
Roots are counted without multiplicity, and only integer roots are counted. Arbitrary integer constants cannot be inserted free of cost; allowing that changes the complexity measure. The degree can grow exponentially in program length, so a degree bound is not the requested polynomial estimate.
\par\smallskip\textit{Formulation and concept references:} \href{https://www.fim.uni-passau.de/fileadmin/dokumente/fakultaeten/fim/lehrstuhl/muller/SmaleProblems1998.pdf}{[S1]}.

\subsection{Short English statement}
Can a polynomial produced by a short constant-free arithmetic program have more than polynomially many distinct integer roots?
\subsection{Sources}
\begin{itemize}
\item[S1]Steve Smale, Mathematical Problems for the Next Century, 1998.
\url{https://www.fim.uni-passau.de/fileadmin/dokumente/fakultaeten/fim/lehrstuhl/muller/SmaleProblems1998.pdf}.
\textit{Formulation source.}
\end{itemize}
\end{document}
